
\documentclass[11pt,a4paper]{scrartcl}

\usepackage[T1]{fontenc}
\usepackage[utf8]{inputenc}
\usepackage[ngerman]{babel}
\usepackage{lmodern}

\usepackage{mathtools}
\usepackage[version=4]{mhchem}
\usepackage{chemfig}
\usepackage{tikz}
\usepackage{caption}

\title{Chemische Formeln und Abbildungen mit LaTeX}
\author{SFN -- LaTeX-Workshop}
\date{\today}

\begin{document}
	
	\maketitle
	\tableofcontents
	
	\section{Chemische Summenformeln}
	
	Wasser (\ce{H2O}), Methan (\ce{CH4})
	und Schwefelsäure (\ce{H2SO4})
	sind Beispiele für chemische Verbindungen.
	
	\section{Reaktionsgleichungen}
	
	\subsection{Verbrennung von Methan}
	
	\begin{equation}
		\ce{CH4 + 2 O2 -> CO2 + 2 H2O}
		\label{eq:methan}
	\end{equation}
	
	\subsection{Ammoniaksynthese}
	
	\begin{equation}
		\ce{N2 + 3 H2 <=> 2 NH3}
	\end{equation}
	
	\subsection{Chemisches Gleichgewicht}
	
	Die Gleichgewichtskonstante ist:
	
	\begin{equation}
		K_c = \frac{
			[\ce{NH3}]^2
		}{
			[\ce{N2}][\ce{H2}]^3
		}
	\end{equation}
	
	\section{Molekülstrukturen}
	
	\subsection{Einfache Verbindungen}
	
	\begin{center}
		\chemfig{H_3C-CH_3}
	\end{center}
	
	\begin{center}
		\chemfig{H_2C=CH_2}
	\end{center}
	
	\subsection{Essigsäure}
	
	\begin{figure}[htbp]
		\centering
		\chemfig{H_3C-C(=[2]O)-OH}
		\caption{Strukturformel von Essigsäure}
		\label{fig:essigsaeure}
	\end{figure}
	
	\subsection{Benzol}
	
	\begin{figure}[htbp]
		\centering
		\chemfig{*6(-=-=-=)}
		\caption{Kekulé-Struktur von Benzol}
		\label{fig:benzol}
	\end{figure}
	
	\section{Reaktionsschema}
	
	Hydrierung von Ethen:
	
	\begin{center}
		\schemestart
		\chemfig{H_2C=CH_2}
		\arrow{->[\ce{H2}][Ni]}
		\chemfig{H_3C-CH_3}
		\schemestop
	\end{center}
	
	\section{Mathematische Gleichungen in der Chemie}
	
	Die Arrhenius-Gleichung beschreibt die
	Temperaturabhängigkeit der
	Reaktionsgeschwindigkeitskonstante:
	
	\begin{equation}
		k = A \exp\left(
		-\frac{E_a}{RT}
		\right)
	\end{equation}
	
\end{document}
